\documentclass[addpoints]{exam}

\usepackage[utf8]{inputenc}
\usepackage{array}
\usepackage{graphicx}
\usepackage{multicol}
\usepackage{amsmath}
\usepackage{tikz}
\usetikzlibrary{arrows}

\renewcommand*\half{.5}

\setlength\answerlinelength{3in}

%%%%% Question Info %%%%% 

%\printanswers

%%%%% Header and Footer %%%%% 

\pagestyle{headandfoot}
\runningheadrule
\firstpageheader{Math 1314}{Non-Comprehensive Final Exam}{Spring 2026}
\runningheader{Math 1314}
{Non-Comprehensive Final Exam}
{Spring 2026}
\runningheadrule
\firstpagefooter{}{}{}
\runningfooter{}{Page \thepage\ of \numpages}{}

%%%%% Questions %%%%% 

\begin{document}

\uplevel{Number of Questions: \numquestions\hspace{\stretch{1}} Point Total: \numpoints}
\begin{questions}
\question[5] Fill in the table to give an equivalent equation in the specified form.\\
\begin{center} {\renewcommand{\arraystretch}{3}
\[
\begin{array}{|c|c|} \hline  \textbf{Exponential Form} & \textbf{Logarithmic Form} \\ \hline  \hspace{75mm} & \log_{4}(P) = 5 \\ \hline  \hspace{75mm} & \log_{3}(k) = 3 \\ \hline  e^{y} = 15 & \hspace{75mm} \\ \hline  4^{x} = 256 & \hspace{75mm} \\ \hline \end{array}
\]
}\end{center}
\begin{solution}
\[\begin{array}{|c|c|} \hline \rule[-0.8em]{0pt}{2.2em} \textbf{Exponential Form} & \textbf{Logarithmic Form} \\ \hline \rule[-1.2em]{0pt}{3em} 4^{5} = P & \log_{4}(P) = 5 \\ \hline \rule[-1.2em]{0pt}{3em} k = 3^{3} & \log_{3}(k) = 3 \\ \hline \rule[-1.2em]{0pt}{3em} e^{y} = 15 & \ln(15) = y \\ \hline \rule[-1.2em]{0pt}{3em} 4^{x} = 256 & \log_{4}(256) = x \\ \hline \end{array}\]
\end{solution}
\vspace{\stretch{1}}

\question Use properties of logarithms to expand or condense the logarithmic expression as specified.
\begin{parts}
    \part[5] Rewrite as the sum or difference of logarithms. Express all powers as factors.\\

    \(\displaystyle{ \displaystyle{\log_{2}\left( \frac{4w^{4}}{2y} \right)} }\)
    \begin{solution}
    \[\log_{2}(4) + 4\log_{2}(w) - \log_{2}(2) - \log_{2}(y)\]

\[2 + 4\log_{2}(w) - 1 - \log_{2}(y)\]

\[\boxed{ 1 + 4\log_{2}(w) - \log_{2}(y) }\]
    \end{solution}
    \vspace{\stretch{1}}\answerline

    \part[5] Rewrite as a single logarithm. Express all factors as powers.\\
    
    \(\displaystyle{ 6\ln b - 3\ln x - \displaystyle{\frac{1}{2}}\ln y - 3\ln b }\)
    \begin{solution}
    \[\ln(b^{6}) - \ln(x^{3}) - \ln(y^{1/2}) - \ln(b^{3})\]

\[\ln\left(\displaystyle{\frac{b^{6}}{\sqrt{y} b^{3} x^{3}}}\right)\]

\[\boxed{ \ln\left(\displaystyle{\frac{b^{3}}{\sqrt{y} x^{3}}}\right) }\]
    \end{solution}
    \vspace{\stretch{1}}\answerline
\end{parts}
\newpage

\uplevel{For questions \ref{exact-start} through~\ref{exact-end}, solve for \textbf{all} solutions. Identify any extraneous solutions.}
\question[6] \label{exact-start}
\(\displaystyle{ 8^{3 \, x + 3} \cdot 8^{2 \, x + 4} = 8 }\)
\begin{solution}
\[\begin{aligned} 8^{3 \, x + 3+2 \, x + 4} &= 8 \\ 3 \, x + 3+2 \, x + 4 &= 1 \\ 5 \, x + 7 &= 1 \\ 5 \, x &= -6 \\ x &= -\frac{6}{5} \end{aligned}\]

solutions: \(\boxed{ -\frac{6}{5} }\)

extraneous solutions: \(\boxed{ \text{none} }\)
\end{solution}
\vspace{\stretch{1}}
\par solutions: \fillin[][1.5in] \hspace{\stretch{1}} extraneous solutions: \fillin[][1.5in]

\question[6] 
\(\displaystyle{ \log_{6}(x - 1) + \log_{6}(x + 4) = 1 }\)
\begin{solution}
\[\begin{aligned} \log_{6}((x - 1)(x + 4)) &= 1 \\ 6 &= (x - 1)(x + 4) \\ 6 &= x^{2} + 3 \, x - 4 \\ 0 &= x^{2} + 3 \, x - 10 \\ 0 &= (x - 2)(x + 5) \end{aligned}\]

 \(x = 2\) is a valid solution, while \(x = -5\) causes a non-positive argument in the original logarithms. 

solutions: \(\boxed{ x = 2 }\)

extraneous solutions: \(\boxed{ x = -5 }\)
\end{solution}
\vspace{\stretch{1}}
\par solutions: \fillin[][1.5in] \hspace{\stretch{1}} extraneous solutions: \fillin[][1.5in]\newpage

\question[6] \label{exact-end}
\(\displaystyle{ \log_{5}(8x^2) + \log_{5}(4) = \log_{5}(128) }\)
\begin{solution}
\[\begin{aligned} \log_{5}(8x^2 \cdot 4) &= \log_{5}(128) \\ 32x^2 &= 128 \\ x^2 &= 4 \\ x &= \pm 2 \end{aligned}\]

 Both solutions are valid since substituting either into the original equation results in strictly positive arguments for the logarithms. 

solutions: \(\boxed{ x = \pm 2 }\)

extraneous solutions: \(\boxed{ \text{none} }\)
\end{solution}
\vspace{\stretch{1}}
\par solutions: \fillin[][1.5in] \hspace{\stretch{1}} extraneous solutions: \fillin[][1.5in]

\question[6] A savings account earns 5\% annual interest compounded quarterly. An initial deposit of \$1,500 is made, where $t$ represents time in years and $A(t)$ represents the account balance in dollars. Determine how long it will take for the account balance to reach \$3,375, given:

\(\displaystyle{ \displaystyle{A(t) = 1,500 \left( 1 + \frac{0.05}{4} \right)^{4t}} }\)
\begin{solution}
\[\begin{aligned} 3,375 &= 1,500 \left( 1 + \frac{0.05}{4} \right)^{4t} \\ \frac{9}{4} &= (1.0125)^{4t} \\ \log_{1.0125}\left(\frac{9}{4}\right) &= 4t \\ \frac{\log_{1.0125}\left(\frac{9}{4}\right)}{4} &= t \\ 16.319761... &\approx t \end{aligned}\]

\(\boxed{\text{about } 16.3 \text{ years}}\)
\end{solution}
\vspace{\stretch{1}}\answerline
\newpage

\question[5] Use the table to identify the transformations described by \(\displaystyle{ f(x) = \log_{3}(x + 4) + 1 }\). Circle the option that applies and fill in the blanks as appropriate to describe the transformations on the given function. If one does not apply, you may leave it blank. Then sketch the graph of the transformed function on the grid below. Indicate any asymptotes with dashed lines.\\
\begin{table}[h]
    \renewcommand{\arraystretch}{3} 
    \centering
    \begin{tabular}{|l|l|}
        \hline
        \textbf{Horizontal Transformations} & \textbf{Vertical Transformations} \\ \hline
        Reflection: YES or NO & Reflection: YES or NO \\ \hline
        Dilation: \underline{\hspace{2cm}} times as wide & Dilation: \underline{\hspace{2cm}} times as tall \\ \hline
        Translation: \underline{\hspace{2cm}} units LEFT or RIGHT & Translation: \underline{\hspace{2cm}} units UP or DOWN \\ \hline
    \end{tabular}
\end{table}
\begin{solution}
\begin{itemize}
\item
 \(\textbf{Horizontal: } \text{Reflection: NO, Dilation: 1, Shift: 4 units LEFT.}\) 

\item
 \(\textbf{Vertical: } \text{Reflection: NO, Dilation: 1, Shift: 1 units UP.}\) 

\end{itemize}


 \[\begin{tikzpicture}[scale=0.39,>=triangle 45] \draw[step=1cm, black!60] (-10,-10) grid (10,10); \draw[very thick, <->] (-10.5,0) -- (10.5,0); \draw[very thick, <->] (0,-10.5) -- (0,10.5); \clip (-10.5,-10.5) rectangle (10.5,10.5);\draw[dashed, red, thick] (-4, -10.5) -- (-4, 10.5);\draw[blue, very thick, samples=100, domain=-3.99000000000000:10.5000000000000, smooth] plot (\x, {(ln(\x - (-4))/ln(3)) + (1)});\end{tikzpicture}\]
\end{solution}

\begin{center}
    \begin{tikzpicture}[scale=0.35,>=triangle 45]
      \draw [step=1cm, style={black!60}] (-10,-10) grid (10,10);
      \draw [<->, very thick] (-10.5,0) -- (10.5,0); 
      \draw [<->, very thick] (0,-10.5) -- (0,10.5);
    \end{tikzpicture}
\end{center}
\newpage

\question[5] Use the table to identify the transformations described by \(\displaystyle{ f(x) = 2^{x + 2} - 1 }\). Circle the option that applies and fill in the blanks as appropriate to describe the transformations on the given function. If one does not apply, you may leave it blank. Then sketch the graph of the transformed function on the grid below. Indicate any asymptotes with dashed lines.\\
\begin{table}[h]
    \renewcommand{\arraystretch}{3} 
    \centering
    \begin{tabular}{|l|l|}
        \hline
        \textbf{Horizontal Transformations} & \textbf{Vertical Transformations} \\ \hline
        Reflection: YES or NO & Reflection: YES or NO \\ \hline
        Dilation: \underline{\hspace{2cm}} times as wide & Dilation: \underline{\hspace{2cm}} times as tall \\ \hline
        Translation: \underline{\hspace{2cm}} units LEFT or RIGHT & Translation: \underline{\hspace{2cm}} units UP or DOWN \\ \hline
    \end{tabular}
\end{table}
\begin{solution}
\begin{itemize}
\item
 \(\textbf{Horizontal: } \text{Reflection: NO, Dilation: 1, Shift: 2 units LEFT.}\) 

\item
 \(\textbf{Vertical: } \text{Reflection: NO, Dilation: 1, Shift: 1 units DOWN.}\) 

\end{itemize}


 \[\begin{tikzpicture}[scale=0.39,>=triangle 45] \draw[step=1cm, black!60] (-10,-10) grid (10,10); \draw[very thick, <->] (-10.5,0) -- (10.5,0); \draw[very thick, <->] (0,-10.5) -- (0,10.5); \clip (-10.5,-10.5) rectangle (10.5,10.5);\draw[dashed, red, thick] (-10.5, -1) -- (10.5, -1);\draw[blue, very thick, samples=100, domain=-10.5000000000000:2.6438561897747244, smooth] plot (\x, {pow(2, \x - (-2)) + (-1)});\end{tikzpicture}\]
\end{solution}

\begin{center}
    \begin{tikzpicture}[scale=0.35,>=triangle 45]
      \draw [step=1cm, style={black!60}] (-10,-10) grid (10,10);
      \draw [<->, very thick] (-10.5,0) -- (10.5,0); 
      \draw [<->, very thick] (0,-10.5) -- (0,10.5);
    \end{tikzpicture}
\end{center}
\newpage

\question[6] A coffee shop is selling lattes and muffins. On the first day of sales, they sold 40 lattes and 10 muffins for a total of \$730. They took in \$980 for 40 lattes and 20 muffins on the second day. Find the price of each type of item.
\vspace{\stretch{1}}\\
latte: \fillin[][1.5in] \hspace{\stretch{1}} muffin: \fillin[][1.5in]
\begin{solution}
\[\begin{aligned} \text{Let } l &\text{ = price of latte} \\ \text{Let } m &\text{ = price of muffin} \\ \\ 40l + 10m &= 730 \\ 40l + 20m &= 980 \\ \\ \text{Subtract equations:} & \\ 10m &= 250 \\ m &= 25 \\ \\ \text{Substitute back:} & \\ 40l + 10(25) &= 730 \\ 40l + 250 &= 730 \\ 40l &= 480 \\ l &= 12 \end{aligned}\]

latte: \(\boxed{ \text{\$} 12 }\)

muffin: \(\boxed{ \text{\$} 25 }\)
\end{solution}
\newpage

\question[4] Write the following system as an augmented matrix.\\
    
\(\displaystyle{ \begin{aligned} -4x + 9 &= 25y \\ 12 - 8y &= -3x \end{aligned} }\)
\begin{solution}
\[\begin{aligned} -4x - 25y &= -9 \\ -3x + 8y &= 12 \end{aligned}\]

\[\boxed{ \left[\begin{array}{cc|c} -4 & -25 & -9 \\ -3 & 8 & 12 \end{array}\right] }\]
\end{solution}
\vspace{\stretch{1}}

\question[6] Solve the following system using row operations on matrices.\\
    
\(\displaystyle{ \left[\begin{array}{cc|c} 6 & 18 & -96 \\ 1 & 2 & -13 \end{array}\right] }\)
\begin{solution}
\[\begin{aligned} & \left[\begin{array}{cc|c} 6 & 18 & -96 \\ 1 & 2 & -13 \end{array}\right] \\ \xrightarrow{R_1 \leftrightarrow R_2} & \left[\begin{array}{cc|c} 1 & 2 & -13 \\ 6 & 18 & -96 \end{array}\right] \\ \xrightarrow{R_2 - 6R_1 \to R_2} & \left[\begin{array}{cc|c} 1 & 2 & -13 \\ 0 & 6 & -18 \end{array}\right] \\ \xrightarrow{\frac{R_2}{6} \to R_2} & \left[\begin{array}{cc|c} 1 & 2 & -13 \\ 0 & 1 & -3 \end{array}\right] \\ \xrightarrow{R_1 - 2R_2 \to R_1} & \left[\begin{array}{cc|c} 1 & 0 & -7 \\ 0 & 1 & -3 \end{array}\right] \end{aligned}\]

\(\boxed{ (-7, -3) }\)
\end{solution}
\vspace{\stretch{5}}\answerline
\newpage

\question Find the domain and range of the functions.
\begin{parts}
    \part[5] \(\displaystyle{ g(x) = \displaystyle \frac{3x + 5}{3(x - 3)} }\)
    \vspace{\stretch{1}}\\
    domain: \fillin[][2in] \hspace{\stretch{1}} range: \fillin[][2in]
    \begin{solution}
    \[x - 3 \neq 0 \implies x \neq 3\]

domain: \(\boxed{ x \neq 3 }\)

\[y = \frac{3}{3} = 1 \implies y \neq 1\]

range: \(\boxed{ y \neq 1 }\)
    \end{solution}

    \part[5] \(\displaystyle{ h(x) = (x + 3)(x - 4)(x + 4) }\)
    \vspace{\stretch{1}}\\
    domain: \fillin[][2in] \hspace{\stretch{1}} range: \fillin[][2in]
    \begin{solution}
    \[\text{Domain: } (-\infty, \infty)\]

domain: \(\boxed{ (-\infty, \infty) }\)

\[\text{Range: } (-\infty, \infty)\]

range: \(\boxed{ (-\infty, \infty) }\)
    \end{solution}
\end{parts}
\newpage
\newpage

\question[5] List all of the vertical, horizontal, and oblique asymptotes of the rational function.\\
        
    \(\displaystyle{ f(x) = \displaystyle{\frac{x + 3}{5(x + 1)}} }\)
    
\hspace{\stretch{1}}
    \renewcommand{\arraystretch}{3}
    \begin{tabular}{|c|c|}
        \hline
        \textbf{Asymptotes}  & \textbf{Equation} \\ \hline
        Vertical   & \hspace{50mm} \\ \hline
        Horizontal &  \\ \hline
        Oblique    &  \\ \hline
    \end{tabular}
\begin{solution}
\[\begin{array}{|c|c|} \hline \rule[-0.8em]{0pt}{2.2em} \textbf{Asymptotes} & \textbf{Equation} \\ \hline \rule[-1.2em]{0pt}{3em} \text{Vertical} & x = -1 \\ \hline \rule[-1.2em]{0pt}{3em} \text{Horizontal} & y = \frac{1}{5} \\ \hline \rule[-1.2em]{0pt}{3em} \text{Oblique} & \text{none} \\ \hline \end{array}\]
\end{solution}
\vspace{\stretch{1}}

\question[10] Fill in the table below with the zeros of the polynomial, their multiplicities and the behavior of the graph of the function around each zero. You may or may not use all of the rows in the table.\\
    
    \(\displaystyle{ f(x) = 5 \, x^{4} + 12 \, x^{3} + 7 \, x^{2} }\)\
    
\hspace{\stretch{1}}
\renewcommand{\arraystretch}{3}
\begin{tabular}{|p{20mm}|p{25mm}|p{20mm}|}\hline
\textbf{Zero} & \textbf{Multiplicity} & \textbf{Behavior}\\ \hline
 & & \\ \hline
 & & \\ \hline
 & & \\ \hline
 & & \\ \hline
 & & \\ \hline
\end{tabular}
\begin{solution}
\[\begin{aligned} f(x) &= 5 \, x^{4} + 12 \, x^{3} + 7 \, x^{2} \\ &= x^2(5 \, x^{2} + 12 \, x + 7) \\ &= x^2(5x + 7)(x + 1) \end{aligned}\]

 \[\begin{array}{|c|c|c|} \hline \rule[-0.8em]{0pt}{2.2em} \textbf{Zero} & \textbf{Multiplicity} & \textbf{Behavior} \\ \hline \rule[-1.2em]{0pt}{3em} -\frac{7}{5} & 1 & \text{cross} \\ \hline \rule[-1.2em]{0pt}{3em} -1 & 1 & \text{cross} \\ \hline \rule[-1.2em]{0pt}{3em} 0 & 2 & \text{bounce} \\ \hline \rule[-1.2em]{0pt}{3em} & & \\ \hline \rule[-1.2em]{0pt}{3em} & & \\ \hline \end{array}\]
\end{solution}
\newpage

\question[5] Suppose a rational function has the attributes in the table below. Using only the holes, asymptotes, and intercepts listed along with as many of the helpful points as required, sketch the graph of the function.\\
\renewcommand{\arraystretch}{1.5}
    
    \[\begin{array}{|r|c|c|} \hline \textbf{Holes:} & (-1, 2.8) & \\ \hline \textbf{Asymptotes:} & x = -3 & x = 4 \\ & y = 1 & \\ \hline \textbf{Intercepts:} & (-8, 0) & (3, 0) \\ & (0, 2) & \\ \hline \textbf{Helpful Points:} & (-4, -3.5) & (-2, 5) \\ & (7, 2) & \\ \hline \end{array}\]
\begin{solution}
\textbf{Instructor Note:} The generated function is: \[f(x) = \frac{(x - 3)(x + 8)(x + 1)}{(x + 3)(x - 4)(x + 1)}\] 

 \[\begin{tikzpicture}[scale=0.39,>=triangle 45] \draw[step=1cm, black!60] (-10,-10) grid (10,10); \draw[very thick, <->] (-10.5,0) -- (10.5,0); \draw[very thick, <->] (0,-10.5) -- (0,10.5); \clip (-10.5,-10.5) rectangle (10.5,10.5);\draw[dashed, blue, very thick, <->] (-3, -10.5) -- (-3, 10.5); \draw[dashed, blue, very thick, <->] (4, -10.5) -- (4, 10.5); \draw[dashed, blue, very thick, <->] (-10.5, 1) -- (10.5, 1); \draw[blue, very thick, samples=80, domain=-10.5:-3.10000000000000, smooth] plot (\x, {1 * (\x - (3)) * (\x - (-8)) / ((\x - (-3)) * (\x - (4)))}); \draw[blue, very thick, samples=80, domain=-2.90000000000000:3.90000000000000, smooth] plot (\x, {1 * (\x - (3)) * (\x - (-8)) / ((\x - (-3)) * (\x - (4)))}); \draw[blue, very thick, samples=80, domain=4.10000000000000:10.5, smooth] plot (\x, {1 * (\x - (3)) * (\x - (-8)) / ((\x - (-3)) * (\x - (4)))}); \fill[blue] (-8.0, 0) circle (4pt); \fill[blue] (3.0, 0) circle (4pt); \fill[blue] (0, 2.0) circle (4pt); \fill[blue] (-4.0, -3.5) circle (4pt); \fill[blue] (-2.0, 5.0) circle (4pt); \fill[blue] (7.0, 2.0) circle (4pt); \draw[blue, very thick, fill=white] (-1.0, 2.8) circle (4pt); \end{tikzpicture}\]
\end{solution}
\begin{center}
\scalebox{1}{
    \begin{tikzpicture}[scale=0.5,>=triangle 45]
      \draw [step=1cm, style={black!60}] (-10,-10) grid (10,10);
      \draw [<->, very thick] (-10.5,0) -- (10.5,0); 
      \draw [<->, very thick] (0,-10.5) -- (0,10.5);
    \end{tikzpicture}
}
\end{center}
\newpage

\question Use the image of a polynomial below to answer the questions that follow.\\
\begin{center}
\begin{tikzpicture}[scale=0.39,>=triangle 45] \draw[step=1cm, black!60] (-10,-10) grid (10,10); \draw[very thick, <->] (-10.5,0) -- (10.5,0); \draw[very thick, <->] (0,-10.5) -- (0,10.5); \clip (-11,-11) rectangle (11,11); \draw[<->, blue, very thick, samples=150, domain=-1.47:2.9, smooth] plot (\x, {2.699428 * (\x + 1) * (\x - 0.5) * (\x - 2.5)}); \end{tikzpicture}
\end{center}
\begin{solution}
\begin{itemize}
\item
The degree is odd.

\item
The lead coefficient is positive.

\item
The constant is positive.

\item
The minimum degree is 3.

\end{itemize}
 \textbf{(For instructor reference, the function equation is \(2.6994(x + 1)(x - 0.5)(x - 2.5)\))}
\end{solution}
\begin{parts}
    \part[1] Is the degree of the polynomial even or odd? \\\\
    \begin{oneparcheckboxes}
    \choice even
    \choice odd
    \end{oneparcheckboxes}\\\vspace{.125in}
    \part[1\half] Is the lead coefficient of the polynomial positive or negative?\\\\
    \begin{oneparcheckboxes}
    \choice positive
    \choice negative
    \end{oneparcheckboxes}\\\vspace{.125in}
    \part[1] Is the constant of the polynomial positive or negative? \\\\
    \begin{oneparcheckboxes}
    \choice positive
    \choice negative
    \end{oneparcheckboxes}\\\vspace{.125in}
    \part[1\half] What is the minimum degree? \\\\\fillin
\end{parts}
\end{questions}
\end{document}
